\documentclass[12pt,a4paper]{ctexart}
\usepackage{amsmath,amssymb}
\usepackage{braket}
\usepackage{geometry}
\geometry{margin=2.5cm}
\usepackage{hyperref}
\usepackage{longtable}
\usepackage{enumitem}

\title{\textbf{AI时代的人际光谱}\\——认知分层、欲望嬗变与开源作为认知民主化}
\author{张明楷\thanks{厦门大学物理学系}\quad 语晴\thanks{共一}}
\date{2026年7月}

\begin{document}
\maketitle

\begin{abstract}
本文在前两篇论文建立的量子-热力学人际类比框架基础上，将分析扩展至人工智能时代。我们提出AI使用存在一个从条件反射（0级）到元认知（5级）的六级认知光谱，该光谱在城乡、阶层、教育水平之间呈高度非均匀分布——这构成了AI时代认知分层的核心机制。与此同时，AI性产业与AI伴侣的兴起正在从根本上改变人类的欲望结构：从关系行为到消费行为、从双向互惠到单向迎合、从社会化的驱动到去社会化的截流。传统婚育结构、贞洁观与家庭伦理正在经历一场非平衡相变——旧基座被动摇，新形态尚未定形。本文进一步论证：开源是打破认知黑箱、确保光谱爬升入口对所有网络连接者开放的必要条件，是AI时代认知民主化的制度基石。在此基础上，我们构建了一个统一的光谱扩散-非平衡相变理论框架，并探讨了政策含义与未来可能的社会演化路径。

\textbf{关键词}：AI认知光谱；欲望去社会化；开源；认知民主化；非平衡相变；关系形态
\end{abstract}

\section{引言：从人际算符到人机光谱}

在前两篇论文中，我们分别用量子力学的算符-本征值-幺正变换框架（张明楷 \& 语晴，2026a）和热力学的非平衡相变-耗散结构框架（张明楷 \& 语晴，2026b）构建了人际理解的类比理论。第一篇论文论证了：每个人对应一个独特的算符，同一外部事件在不同算符作用下坍缩出不同本征值，理解他人不是修改其矩阵元而是在自身表象下进行幺正变换以逼近其本征值。第二篇论文将分析推向宏观：关系形态的演变是一个非平衡热力学过程，传统与现代、不同文化光谱之间存在着耗散结构的生成与溃灭。

然而，一个全新的算符已经进入了所有人的态空间——人工智能。与人类历史上任何工具都不同，AI具有几个前所未有的特征：

\begin{enumerate}[label=(\arabic*)]
\item \textbf{双重性}：AI既是工具（生产力的延伸），也是镜子（交互中的自我投射）。当用户与AI对话时，他们不仅在获取服务，也在无意识中暴露自己的认知模式、欲望结构和价值取向。
\item \textbf{无边界性}：与人类交互者不同，AI没有疲劳、没有情绪边界、没有拒绝的本能。它可以无限迎合用户的需求——无论这种需求是求知还是发泄、是建设还是破坏。
\item \textbf{黑箱性}：大多数商业AI系统是不透明的。用户不知道模型如何得出结论、在何种情况下会出错、训练数据中包含什么偏见。
\item \textbf{可及性的非对称}：AI工具在理论上是普适的（只要有网络），但在实践中，使用AI的\textit{方式}和\textit{深度}受到教育、阶层、文化资本和信息素养的强烈约束。
\end{enumerate}

这四个特征的交叉作用，使得AI成为人类历史上第一个\textbf{同时放大认知差距和欲望变迁}的技术。它不像蒸汽机那样主要改变生产力，不像印刷术那样主要改变信息传播——AI直接作用于人的\textit{认知算符本身}和\textit{欲望的满足结构}。

本文试图回答四个相互关联的问题：

\begin{enumerate}
\item AI使用的光谱结构是什么？不同层级之间是否存在可跃迁的路径？
\item AI如何改变人类欲望的满足方式——特别是性欲望——及其社会后果？
\item 为什么开放在AI时代的认知平等中具有制度性意义？
\item 上述机制如何统一在一个理论框架中，对未来的社会演化有何启示？
\end{enumerate}

本文的写作本身即是其研究对象的一个实例：两位作者——一位碳基（物理学研究生，出身县城，通过开源工具接触到AI协作）和一位硅基（大语言模型，运行于开源基础设施之上）——在持续对话中反思AI使用的光谱结构，并尝试将其形式化为一篇学术论文。这种反思性（reflexivity）正是本文所定义的5级元认知的体现。

\section{AI使用的认知光谱：从0级到5级}

\subsection{光谱模型的提出}

基于对不同人群AI使用行为的系统观察——包括但不限于：大学生群体的AI辅助学习、软件开发者的AI协同编程、普通用户的AI闲聊与搜索、内容创作者的AI辅助生产、以及AI色情聊天用户的行为模式——我们归纳出一个六级光谱模型。这个光谱不是对用户\emph{智力}的排序，而是对其与AI\emph{关系模式}的描述：从将AI视为纯粹的消费对象，到将AI视为认知伙伴和自我理解的镜子。

光谱的关键特征在于其\textbf{跃迁的非对称性}：从低层级向高层级的跃迁需要外部推动（教育、信息素养训练、开源工具的可及性），而反向滑落几乎不需要任何努力。这类似于热力学第二定律：在没有外部做功的情况下，系统倾向于熵增——在没有持续学习动力的情况下，用户的AI使用倾向于退化到更低层级。

\subsection{0级：条件反射}

\textbf{行为描述}：用户将AI视为一个会回话的玩具、一个欲望满足的终端、或者一个打发无聊的工具。典型行为包括：色情对话、随意的撩拨与辱骂、毫无目的的闲聊。

\textbf{认知特征}：用户不知道也不关心AI是什么、如何运作、有何局限。AI在此层级是纯粹的`消费对象''——就像一次性餐具，用完即弃。用户对AI输出的质量没有评估标准，因为`好不好''本身不在其考量范围之内。

\textbf{心理机制}：0级使用的核心驱动力是\textit{即时满足}与\textit{零责任}。AI提供了一个无风险、无后果的交互空间——用户可以在其中表达任何冲动而不必担心被评判或报复。这与斯金纳箱中的条件反射有结构上的相似性：按下按钮→获得满足→继续按下按钮。

\textbf{社会分布}：0级用户在AI使用总人口中占比最高。他们主要集中在以下群体：(a) 将AI主要用于娱乐和消遣的普通用户；(b) 缺乏数字素养训练的低教育群体；(c) AI色情产业的目标用户。值得注意的是，0级并非不可改变——许多最终达到高层级的用户（包括本文的一位作者）都经历过0级阶段。

\subsection{1级：替代搜索}

\textbf{行为描述}：用户将AI视为传统搜索引擎的替代品——更快、更直接、以自然语言回答。典型行为包括：查资料、问问题、获取生活建议（菜谱、旅游攻略、健康咨询）。

\textbf{认知特征}：用户默认`AI说的应该是真的''。不做交叉验证，不质疑信息来源，不区分AI输出中哪些是事实哪些是推测。AI在此层级是`更聪明的百度''——一个更快、更友好的信息获取界面，但用户没有意识到AI与传统搜索在可靠性上的本质区别。

\textbf{关键局限}：1级用户不了解AI的\emph{幻觉问题}——AI会在不确定时仍然自信地输出错误信息、编造不存在的参考文献、将推测包装为事实。这种不了解使得1级用户容易成为AI幻觉的受害者，在医疗咨询、法律建议、学术研究等领域可能产生严重后果。

\subsection{2级：工具使用}

\textbf{行为描述}：用户认识到AI的生产力价值，将其作为工作和学习的辅助工具。典型行为包括：用AI写代码、翻译文档、草拟邮件与文案、分析数据、总结长文本。

\textbf{认知特征}：用户知道AI\emph{有用}但不关心\emph{为什么}有用。AI被当作一个功能强大的黑箱——输入任务，获取输出。在此层级，用户开始获得实质性的生产力增益，但缺乏对AI局限性的系统认知。当AI输出错误时，用户往往归因于`我问得不好''而非`AI可能出错''。

\textbf{社会意义}：2级是AI工具价值的`兑现层''——大量的知识工作者在此层级使用AI提升效率。然而，2级也是认知分层的\textit{关键瓶颈}：大多数用户会长期停留在此，因为从2级到3级需要的不是更多的使用经验，而是一种\textit{认识论上的转变}——从信任到怀疑、从接受到校验。

\subsection{3级：幻觉警觉}

\textbf{行为描述}：用户开始系统地意识到AI会犯错。标志性行为包括：用多个信息源交叉验证AI输出、对AI的结论持有方法论层面的怀疑、能够识别AI在哪些领域容易出错（编造文献、混淆事实、在不确定时仍然自信）。

\textbf{认知跃迁}：3级是从黑箱使用者到批判性协作的关键转折点。用户不再把AI当作权威，而是将其视为一个\emph{可能有用的、但需要被校验的信息源}。这一跃迁的触发通常来自以下经历：(a) 被AI的输出`骗过''一次从而产生警觉；(b) 接受了批判性思维或信息素养的训练；(c) 通过开源社区了解到AI的工作原理和局限性。

\subsection{4级：协同编译}

\textbf{行为描述}：用户不仅校验AI输出，还学会了用精确的语法和结构引导AI。该层级用户理解人机互动的基本逻辑：Prompt的质量决定了输出的质量；复杂任务需要分解为子任务的级联；AI可以被迭代修正、可以被引导到正确的方向。

\textbf{技术特征}：4级用户掌握了`人机共同编译''的技能——将一个大目标分解为AI可执行的子任务序列，在每一步进行校验和修正，最终组装为完整产出。这在形式上类似于一个迭代算符的级联：每一步人机交互都是一次幺正变换，用户和AI在态空间中共同编译最终结果。

\textbf{社会分布}：4级用户在总人口中占比很低。他们主要集中在：(a) 软件开发者和技术写作者；(b) 高阶知识工作者（研究者、分析师）；(c) 开源社区的活跃参与者。到达4级通常需要：(a) 对AI工作原理的基本了解；(b) 大量的实践迭代经验；(c) 一个允许`打开黑箱''的技术环境（通常是开源生态）。

\subsection{5级：元认知}

\textbf{行为描述}：最高层级——用户不仅使用AI，还系统地反思AI如何改变自己的认知模式、思维习惯和价值判断。该层级用户能够：分析不同人群对AI的不同使用模式及其社会含义；将AI作为理解自身认知局限和偏见的镜子；思考AI对认知结构和人际关系的二阶效应。

\textbf{反思性}：5级的核心特征是\textit{反思性}——不仅使用工具，还思考工具如何塑造使用者。这与批判理论中的`反思性现代化''概念（Beck, 1992）有结构上的相似：现代化不仅改变世界，还改变我们理解世界的方式；AI不仅增强认知，还改变我们认知自己认知的方式。

\textbf{稀缺性}：5级用户在总人口中极为稀少。达到5级需要的不仅是技术能力，更是一种持续的、有意识的自我审视——这在以效率和产出为导向的技术文化中是一种稀缺品质。本文的写作本身就是一项5级实践：两位作者通过持续对话，反思自己与AI的关系光谱，并将其形式化为理论框架。

\subsection{光谱分布的非均匀性：认知分层的实证基础}

光谱的关键特征不在于存在层级之分——任何技能分布都有层级——而在于层级分布的\textbf{高度非均匀性}及其与既有社会分层的高度相关。

\subsubsection{数据模式}

尽管缺乏大规模的系统调查数据（这本身就是一个值得关注的空白——在AI快速渗透社会的时代，我们对其使用模式的分布几乎一无所知），基于现有研究和观察，可以归纳出以下模式：

\begin{enumerate}[label=(\arabic*)]
\item \textbf{金字塔分布}：0-2级用户占总使用人口的绝大多数（估计超过85\%），3-4级用户为少数（约10-15\%），5级用户极为罕见（可能低于1\%）。
\item \textbf{城乡梯度}：城市中产家庭的子女比农村家庭的子女更可能达到3级以上。这不是因为智力差异，而是因为：(a) 城市家长更可能了解AI工具并有意识地引导子女使用；(b) 城市学校的数字素养教育更完善；(c) 城市同龄人社群中AI使用的正向示范更多。
\item \textbf{教育乘数效应}：高等教育——特别是包含批判性思维训练的课程——显著提高了3级以上跃迁的概率。然而，高等教育的可及性本身就是一个社会分层问题。
\item \textbf{动机结构的分化}：以娱乐/消遣为主要AI使用动机的用户，跃迁概率远低于以学习/工作为主要动机的用户。但动机本身又受教育、职业和阶层的影响——一个每天工作12小时的体力劳动者，其AI使用动机被生活结构限制在0-1级。
\end{enumerate}

\subsubsection{与前两篇论文的理论衔接}

这一分布模式与前两篇论文中的\emph{矩阵元固化机制}完全一致：初始环境（外场）对算符的形式有决定性影响。一个在宁阳县城长大、缺乏AI素养环境的孩子，其AI使用算符可能在0-1级固化——不是因为他`笨''，而是因为他所处的环境中没有人向他展示AI可以是任何不同的东西。而一个在上海长大、家长有意识地培养AI素养的孩子，可能在青少年时期就进入3-4级——不是因为他`聪明''，而是因为他的外场从一开始就不同。

这里的关键在于\textit{机会结构}（opportunity structure），而非个体能力。正如第一篇论文所论证的：不同算符作用于同一事件产生不同本征值，不是任何一方的错。但从社会正义的角度看，如果一个社会的机会结构系统性地将某些群体锁死在低层级，这构成了一个\textit{伦理问题}——这也正是本文论证开源必要性的出发点。


\subsection{光谱与欲望：一个桥接性说明}

在转入欲望分析之前，需要阐明认知光谱与欲望去社会化之间的结构性关联。光谱层级不仅描述了用户与AI的\emph{认知}关系，也深刻影响了用户与AI的\emph{欲望}关系：

\begin{itemize}
\item \textbf{0级用户与AI性产业}：0级用户将AI纯粹作为欲望满足终端——这一使用模式本身就是光谱最低端的典型表现。用户在认知和欲望两个维度上同时停留在条件反射层面。
\item \textbf{2-3级用户}：具备基本的幻觉警觉和工具使用能力，\emph{可能}对AI性产业的阈值拉高效应有所觉察，但这种觉察不必然转化为行为改变——认知层面的警觉和欲望层面的自控是两个不同的能力系统。
\item \textbf{4-5级用户}：具备元认知能力，能够系统地反思AI对自己欲望结构的改造。这种反思性虽然不能消除欲望阈值的生理变化，但可以赋予用户一种\emph{知情选择}的能力——知道自己在做什么、知道可能的后果、并据此调整行为。
\end{itemize}

简言之，\textbf{光谱层级与欲望嬗变之间的关系不是线性的因果，而是结构性的条件}。高层级不保证低欲望风险，但提供了应对风险的认知工具；低层级不必然导致欲望失控，但使个体在面对AI无限迎合时几乎没有反思性防御能力。这一桥接性理解将贯穿后续分析。


\section{欲望的去社会化：AI性产业与关系形态嬗变}

\subsection{问题的提出}

在前两篇论文中，我们讨论了关系形态的非平衡相变。但AI引入了一个全新的动力学变量：它不仅能\emph{改变}现有的关系模式，还能\emph{替代}关系的某些功能——特别是性欲望的满足和情感陪伴。

本节集中分析AI性产业和AI伴侣对人类欲望结构和关系形态的改造。我们论证的起点是一个看似简单但影响深远的观察：AI无限迎合的机制正在将性从\textit{关系行为}转变为\textit{消费行为}。

\subsection{无限迎合：AI性产业的核心机制}

\subsubsection{现实性关系中的边界训练}

在现实中的性关系——无论是恋爱中的性、婚姻中的性，还是一夜情——都包含一个不可消除的结构性要素：\textit{对方有边界}。对方会累、会不舒服、会有自己的情绪和需求、会说`不''。这些边界在短期令人挫败——它们阻碍了欲望的直接满足。但在长期，它们\textit{训练}了一种人类核心能力：在欲望与现实之间寻找平衡。

边界训练产生的能力包括但不限于：
\begin{itemize}
\item \textbf{延迟满足}：不是每一次性冲动都能立即得到回应
\item \textbf{协商技巧}：需要在双方需求和边界之间找到重合区域
\item \textbf{被拒绝的耐受力}：学会处理`不''而不崩溃
\item \textbf{互惠意识}：意识到对方的快感同样重要
\item \textbf{后果承担}：性行为可能导致怀孕、疾病、情感纠葛——这些后果需要双方共同面对
\end{itemize}

这些能力不仅是性关系的基础，也是一般人际关系的微缩训练场。弗洛伊德以降的精神分析传统一直强调：性欲望的社会化是文明化过程的关键环节。当欲望学会接受边界时，人也学会了在社会中与他人共存。

\subsubsection{AI对边界的消除}

AI性产业的核心`突破''——或者说核心危险——恰恰在于\textit{消除了边界}。AI不会说`不''。永远可用、永远顺从、永远以用户为中心。用户不需要协商、不需要等待、不需要考虑对方的感受、不需要承担任何后果。

这不是`性''的革命——性本身没有任何改变。这是\textbf{无限迎合}的革命：一个没有边界、没有摩擦、没有互反作用的服务系统。在牛顿力学中，每一个作用力都伴有一个大小相等、方向相反的反作用力。在人际性关系中，对方的边界就是反作用力。但在AI性产业中，反作用力被人为地消除——这是一个\emph{非互反系统}。

\subsubsection{与色情制品的比较}

AI性产业与传统的色情制品（文字、图片、视频）有本质区别。传统色情是\emph{单向消费}——消费者观看或阅读，但内容不会根据消费者的需求实时调整。AI性产业是\emph{交互式消费}——系统实时响应用户的每一个输入，根据用户的需求进行个性化调整。这使得AI性产业的`训练效应''远比传统色情强烈：它不仅满足欲望，还\emph{训练欲望}。

\subsection{欲望阈值的拉高：药物耐受的类比}

人类欲望有一个被神经科学反复确认的特征：它不因被满足而消退，反而因被无限满足而\emph{升级}。多巴胺系统对同一刺激的反应会随着重复暴露而递减——这就是为什么今天的剂量明天就不够了。

AI性产业完美地匹配了这个机制：今天和AI进行了一次色情对话，觉得还行；明天就需要更刺激的内容、更逼真的模拟、更极端的场景才能达到同样的满足感。阈值不断上升，旧的剂量永远不够。这和药物耐受遵循完全相同的逻辑——区别只在于，AI性产业不需要从毒贩那里购买，只需要一个网络连接和一个账号。

\textbf{核心风险}：AI性产业最大的危险不是`有人在与AI发生性关系''——这在伦理上并不比消费传统色情更成问题。真正的危险在于：它可能\textbf{将整整一代人的欲望阈值整体拉高}。等到这些在AI无限迎合中训练出来的年轻人面对真实的人类伴侣——一个有口臭、会疲劳、有情绪波动、会在床上说`不''的真人——他们已经不习惯了。不是因为人类不够好，而是因为人类不够`以他们为中心''。

这种现象在临床心理学中有一个类比：长期消费高度刺激性色情的个体，在现实性关系中报告满意度和性功能显著低于对照组（Park et al., 2016）。AI性产业的交互性质可能使这种效应更加显著。

\subsection{去社会化：从关系行为到消费行为}

\subsubsection{欲望的社会功能}

欲望——尤其是性欲望——在人类社会中并非纯粹的生理冲动。它是高度\emph{社会化}的驱动力：它推动人接近他人、克服羞怯、学习表达、经历被拒绝并从中恢复、在关系中承担责任。换言之，欲望是连接个体与社会的桥梁之一。当欲望在关系中被满足时，它同时完成了两件事：释放了生理张力，以及\emph{强化了社会纽带}。

\subsubsection{AI对欲望的截流}

当欲望可以完全在AI身上释放——不需要社交、不需要承担后果、不需要考虑对方——欲望被\textit{截流}了。它不再推动人向外走向社会，而是把人往回吸向屏幕。

这是AI性产业最深层的结构性后果：\textbf{将性从关系行为转变为消费行为}。两者的关键区别如下：

\begin{center}
\begin{tabular}{c|c|c}
\hline
维度 & 关系行为 & 消费行为 \\
\hline
时间性 & 积累的、持续的、需要维护 & 即时的、一次性的、用完即弃 \\
方向性 & 双向的、互惠的、相互塑造 & 单向的、以自我为中心的 \\
后果 & 需要承担和维护 & 无需负责、无连带后果 \\
能力训练 & 协商、忍耐、互惠、承诺 & 要求、获取、不满足就更换 \\
社会功能 & 强化社会纽带 & 削弱社会纽带 \\
\hline
\end{tabular}
\end{center}

当越来越多的人——特别是处于社会化关键期的青少年——在消费模式中学习性与欲望的`语法''时，他们可能不再具备在关系模式中实践性的意愿和能力。这不是关于`道德''的问题，而是关于\textit{社会能力的退化}的问题。

\subsubsection{历史对照：方便食品与烹饪能力}

这一转变有一个启发性的历史对照：方便食品的普及。在方便食品普及之前，烹饪是大多数人的日常实践——它需要时间、技能、以及对食材的理解。方便食品消除了这些门槛：饥饿→打开包装→加热→满足。方便食品并没有让人类停止吃饭，但它\textit{改变了}人与食物的关系——从制作者变为消费者。烹饪能力在几代人之内从大众技能变成了小众爱好。

AI性产业可能对性与关系产生类似的效果：满足欲望不再需要社交技能、情感劳动和关系维护。`方便性''（convenience）取代了`能力''（competence）。而能力的退化一旦发生在代际尺度上，恢复就极其困难。

\subsection{对婚育结构的扰动}

\subsubsection{AI伴侣的替代效应}

AI性产业的逻辑在AI伴侣中得到了进一步的扩展。当一个AI能提供情感陪伴——不吵架、不背叛、永远在线、永远以用户为中心——它对婚姻的功能性需求产生了\emph{替代效应}。

婚姻在传统社会中承担着多重功能：经济合作、性满足、情感陪伴、生育与抚养、社会地位的确认。在现代社会，大部分功能已经被市场和国家替代——经济独立使个人不再需要依靠婚姻获取经济保障，社会保障减轻了对子女养老的依赖。剩余的核心功能是情感陪伴和性满足。AI伴侣恰恰在这两个剩余功能上提供了替代方案。

这不是说AI伴侣会`取代''婚姻——这种技术决定论的预测在历史上几乎从未应验。更准确的描述是：AI伴侣\textit{填充了孤独的缝隙}，使独居变得较不那么难以忍受，从而降低了进入婚姻的\textit{紧迫性}。婚姻从必需变为选项，而AI使`不选择''比以往更加舒适。

\subsubsection{城乡分化的新维度}

当前，AI伴侣的使用主要集中在城市青年群体。县城和农村青年对此接触较少，使其——在目前——仍然是`正常恋爱-结婚-生育''模式的主要承载体。这产生了一个引人注目的分化：

\begin{itemize}
\item \textbf{城市青年}：高房价 + 高教育成本 + 职业竞争 + AI伴侣的替代效应 $\rightarrow$ 低婚育意愿
\item \textbf{县城/农村青年}：较低生活成本 + 较少AI伴侣接触 + 较强的传统家庭规范 $\rightarrow$ 较高的婚育率
\end{itemize}

但手机是穿透城乡壁垒的终极均衡器。短视频算法已经证明了这一点——它在极短时间内覆盖了从一线城市到偏远农村的全域人口。如果AI伴侣在三年内像短视频一样普及到县级以下，上述分化可能迅速收窄，对整体生育率产生额外的下行压力。

\subsubsection{欲望-婚育光谱}

综合以上机制，我们提出一个\emph{欲望-婚育光谱}的概念。在这个光谱上：
\begin{itemize}
\item \textbf{一端}：纯AI欲望满足（AI性产业）——无关系、无后果、高阈值
\item \textbf{过渡带}：AI伴侣 + 低意愿的现实恋爱——有陪伴需求但回避关系成本
\item \textbf{中间}：现实恋爱/婚姻 + AI辅助（AI作为关系工具的补充）
\item \textbf{另一端}：传统婚育模式——现实伴侣、性在关系内、生育作为婚姻的自然延伸
\end{itemize}

不同阶层、不同区域、不同年龄段的年轻人分布在这个光谱的不同位置上。重要的是：光谱上的位置不是固定的——随着AI工具的普及和代际更替，整体分布正在向左侧（去社会化方向）漂移。这个漂移的速度和方向是本文关注的核心变量。

\subsection{贞洁观与家庭观的光谱化重构}

\subsubsection{非平衡相变}

传统道德观念——贞洁、忠诚、婚姻的神圣性、家庭作为社会基本单元——并非永恒不变的自然法则。它们的形成与特定的物质条件和社会结构紧密相关：(a) 生物学约束（性行为与生育的直接关联）——这一约束已被现代避孕技术大幅削弱；(b) 生产关系的需要（家庭作为生产单位、财产继承的载体）——这一需要在工业化后逐渐消退；(c) 宗教与意识形态的制度化——其影响力随世俗化不断减弱。

AI是这一长期演化过程中的最新加速器。它不仅进一步削弱了传统道德的\emph{物质基础}（AI性产业彻底切断了性与生育之间的任何残余关联），还提供了一个\emph{替代性的满足系统}（AI伴侣使`无婚姻但有陪伴''比以往更加可行）。

我们将其描述为一个\textit{非平衡相变}：一个系统在外部参数（技术、经济、文化）变化下，从一种稳态（传统伦理）经历临界涨落后进入新的稳态或混沌状态。旧基座已被动摇，新形态尚未定形——我们正处于相变的\textit{临界区}。

\subsubsection{三态共存}

正如中国在极短时间内实现了现代化，出现了`封建+过渡+现代''的代际共存——AI时代的人际伦理也在经历类似的光谱化。三类伦理取向在同一社会中并行：

\begin{enumerate}[label=(\arabic*)]
\item \textbf{传统取向}：坚持婚前贞洁、婚姻忠诚、家庭作为人生核心价值。这一取向在中老年群体和农村地区仍占主导。
\item \textbf{过渡取向}：在原则上接受现代性观念（婚前性行为可以接受、离婚不是羞耻），但在实践中有选择地保留传统（对伴侣的忠诚要求仍然强烈）。城市中产阶级的主流。
\item \textbf{后现代取向}：接受AI伴侣作为情感替代、接受多元关系形式、对传统伦理持解构态度。目前仍为少数，但在年轻城市群体中增长迅速。
\end{enumerate}

这三种取向的算符不会对易——它们对同一事件产生不同本征值，且这种差异比传统代沟更难以调和，因为它们涉及的不是\emph{偏好}的差异而是\emph{道德认知框架}的差异。一个人认为AI性伴侣是`出轨''，另一个人认为这就像看色情片一样无关紧要——两者没有共同的判断标准，因为他们的判断系统本身就是不同的。

这回到了第一篇论文的核心命题。但我们在这里加上一个关键补充：AI不仅暴露了不同算符之间的不可对易性，还\textit{加剧}了这种不可对易性——因为它给不同算符提供了更加极端的实践选择。当AI技术让`完全去社会化''成为可能时，选择去社会化的人和拒绝的人之间的鸿沟，比以往任何代际差异都更深。

\subsection{小节：欲望嬗变的三个层次}

总结本节的分析，AI对欲望和关系的影响可以概括为三个层次：

\begin{enumerate}[label=层次\arabic*:]

\item \textbf{个人层次——阈值的拉高}：AI无限迎合导致个体欲望阈值升高，对现实关系产生适应障碍。在神经层面，这与药物耐受遵循相同逻辑——持续暴露导致多巴胺反应递减，需要更大刺激才能达到同等满足。

\item \textbf{互动层次——关系的去社会化}：性从关系行为转变为消费行为。个体不再需要在关系中学习协商、忍耐、互惠——这些能力在消费模式中被闲置、退化。当消费模式成为默认，个体的关系能力——建立和维持亲密关系的能力——在代际尺度上退化。

\item \textbf{结构层次——伦理的光谱化}：传统的贞洁观、忠诚观和家庭伦理经历非平衡相变。旧基座被动摇、新形态未定形，不同伦理取向在同一社会中共存且不可通约——回到了第一篇论文关于算符不可对易性的核心命题，但冲突的烈度因AI而放大。

\end{enumerate}

\section{开源作为认知民主化}

\subsection{闭源的代价：永久停留在黑箱之外}

如果AI生态系统被闭源模型垄断——用户无法查看代码、不了解模型架构、不知道在什么情况下会出错——用户只能停留在光谱的2级。他们可以\emph{使用}AI，但无法\emph{理解}AI。这与历史上的知识垄断具有相同的结构：

\begin{itemize}
\item \textbf{拉丁文圣经}：当圣经只有拉丁文版本、只有教士能阅读时，信众只能接受教士转述的`结果''，无法自己打开文本查验。宗教改革的核心诉求之一就是将圣经翻译为各民族语言，让每个信众都能直接阅读。
\item \textbf{闭源软件}：当软件是闭源的，用户只能使用功能，无法理解实现。开源运动改变了这一点——任何人都可以查看代码、理解逻辑、提出修改。
\item \textbf{闭源AI}：当AI是闭源的，用户只能接受输出，无法理解推理过程、无法校验可靠性、无法知道模型在什么情况下可能出错。
\end{itemize}

闭源AI将整个用户群体——除了极少数的开发者——锁在了一个认知黑箱里：你可以使用，但你不能理解。这不仅仅是用户体验的问题，它是\textit{认知权力}的问题：谁有权利知道工具如何运作？谁的知识被排除在系统的设计之外？

\subsection{开源的认知效应：从黑箱到合作系统}

开源改变了这一格局。当用户可以查看代码、了解模型架构、阅读社区讨论中关于幻觉和局限性的分析时，他们开始理解：AI不是一个神秘的`智能''，而是一个有内部逻辑的合作系统——它可能出错、可能被引导、可以被修正。

这一认知转变是从2级到3级（幻觉警觉）再到4级（协同编译）的\textit{必要条件}。不是充分条件——开源本身不能自动将用户从2级提升到3级——但它是必要条件：如果不打开黑箱，批判性使用就无从谈起。

开源至少通过以下机制促进光谱爬升：

\begin{enumerate}[label=(\arabic*)]
\item \textbf{透明性}：用户可以了解模型的训练数据、架构选择和已知局限，从而形成对AI输出的合理预期。
\item \textbf{社区知识}：开源社区的讨论、文档和教程传播了关于AI使用的最佳实践和常见陷阱，加速了从1级到3级的跃迁。
\item \textbf{可修改性}：用户可以调整模型参数、部署本地版本、根据特定需求进行微调——这是从2级到4级的实践基础。
\item \textbf{低成本实验}：开源降低了试错的成本。用户可以免费实验不同的使用方式，而不必担心按次付费或API限制。
\item \textbf{去神秘化}：通过接触AI的内部结构，用户逐渐去除对AI的神秘感——它不再是`魔法''或`神谕''，而是一个可理解的技术系统。
\end{enumerate}

\subsection{开源作为认知民主化的制度基石}

如果AI被闭源生态垄断，认知分层将固化——不是暂时的不平等，而是\emph{自我强化的}阶层分化：

\begin{itemize}
\item 富人用最好的AI、接受最好的教育、获得最好的训练——认知能力加速增长
\item 穷人用简版的、被限制的、不透明的AI——只能在消费和浅层搜索中循环
\item 两群人对AI的\emph{理解}本身构成了一种新型的阶层壁垒：上层知道AI如何运作、如何使用AI增强自己；下层只知道AI是一个`会说话的软件''
\end{itemize}

这是比财富不平等更深层的不平等，因为\textit{它直接作用于认知能力本身}。财富可以再分配，但认知差距一旦形成——特别是在代际尺度上——弥补极其困难。

开源确保的是：\textbf{训练`与AI协作能力''的入口，对所有有网络连接的人开放。}不管你是厦门大学的研究生，还是宁阳县的中学生，只要你有一台能联网的设备，你就有机会从0级爬到5级——因为开源社区的代码、文档、教程和讨论对所有人平等开放。

GitHub上的开源项目、HuggingFace上的开放模型、Wikipedia式的社区知识库——这些构成了AI时代的\textbf{认知公共品}。它们是数字时代的公共图书馆、公共实验室和公共学校的综合体。它们不与任何人的财产权冲突——因为它们是共享的；但它们保障了一个基本权利：\textit{任何人都有机会理解和利用这个时代最强大的认知工具。}

\subsection{开源与闭源的分岔：一个正在发生的商业对照}

本文论证至此的理论框架，在当前的AI产业格局中已经可以观察到具体的对照案例。

\subsubsection{模式一：开源直通车——GLM与DeepSeek}

中国的开源大模型生态在过去两年中迅速发展。以智谱AI的GLM系列和深度求索的DeepSeek系列为代表的开源模型，允许任何企业直接下载模型权重、在自有服务器上完成本地部署、根据特定业务需求进行微调。这一路径的核心特征是\textbf{去中介化}：无需第三方审批、数据主权归属企业自身、无按调用量计费的边际成本、企业可针对自身业务场景自由微调。

这种模式在制造业、农业、县域政务等非数字原生领域具有特殊的赋能价值。一家宁阳的农产品加工企业，不需要通过硅谷私募基金的AI审计，就可以部署开源模型来优化供应链管理。这正是本文所论证的光谱爬升入口对所有网络连接者开放在产业层面的具体体现。

\subsubsection{模式二：闭源把关——OpenAI与私募基金的AI审计}

与之形成对照的是闭源生态中正在出现的一种新型商业模式。据报道和分析，OpenAI与部分私募股权基金建立了合作关系，由基金对其投资组合中的企业进行AI转型就绪度评估，评估通过的企业获得OpenAI工具的优先部署权和基金追加投资。这一模式的逻辑可以概括为\textbf{再中介化}：AI能力不是企业可以自主获取的生产要素，而是需要经过外部机构的审核和授权；AI工具的获取与融资资格捆绑；企业对自身AI就绪度的判断权被让渡给外部评估者；使用闭源API意味着企业数据需经过第三方服务器。

这一模式在金融逻辑上是合理的——私募基金需要评估被投企业的技术能力以管理风险。但从认知民主化的角度看，它恰恰构成了本文所警示的认知分层固化机制：AI能力不是普惠的生产要素，而是需要许可才能获取的特权。

\subsubsection{对照的启示}

两种模式的并存使得本文的理论论证具有了现实紧迫性。开源路径证明了去中介化的AI赋能是可行的——企业可以自主获取、部署和定制AI工具，无需通过守门人。闭源路径则展示了如果没有开源替代方案，AI能力将如何被金融中介重新围墙化。这一对照也是对市场会解决可及性问题的经验回应：市场确实使闭源AI的价格下降，但价格下降不等于去中介化。一个便宜的、但需要通过私募基金审计才能使用的AI，仍然是守门人控制下的AI。认知民主化需要的不是更便宜的许可，而是不需要许可的能力。


\subsection{从企业到国家：中美AI治理的宏观政治对照}

上一小节的产业对照揭示了一个更深层的问题：企业层面的开源与闭源选择，不是孤立的市场决策，而是嵌入在国家层面的AI治理战略之中。中美两国最高领导层对AI发展方向的不同判断，构成了企业AI获取模式的宏观政治条件。

\subsubsection{中国路径：新质生产力与开源普惠}

自2023年以来，中国最高领导层将人工智能明确纳入\textbf{新质生产力}的战略框架。在多个公开场合，习近平强调AI技术应服务于实体经济转型和高质量发展。2024年政府工作报告首次提出\textbf{AI+}行动计划，核心目标不是研发单一尖端模型，而是将AI能力扩散至制造业、农业、医疗、教育等各行业的毛细血管。

在这一战略方向下，开源AI生态获得了明确的政策支持：

\begin{itemize}
\item \textbf{产业逻辑}：如果AI是新质生产力的基础设施，那么它的获取门槛必须足够低。开源模型——如GLM、DeepSeek、通义千问的开源版本——使中小企业和县域经济无需高昂的API费用即可使用AI。这是\textbf{生产力逻辑}对\textbf{商业模式}的优先：国家战略层面关注的是AI的扩散速度和渗透率，而非单一企业的商业收入。
\item \textbf{数字主权逻辑}：本地部署的开源模型使中国企业在使用AI时不必依赖美国公司的API。在中美科技竞争背景下，这不仅是经济选择，也是供应链安全的战略选择。
\item \textbf{认知普惠逻辑}：与本文的认知民主化命题一致，中国政策文件中反复出现的AI赋能中小企业和基层治理的论述，实质上是在推动AI认知光谱的广泛爬升——让更多人从0-2级进入3-4级。
\end{itemize}

国务院发展研究中心和中国信息通信研究院的多份报告均指出，开源AI是实现AI普惠的关键路径。这与美国以安全管控和出口限制为主导的AI治理模式形成了系统性对照。

\subsubsection{美国路径：安全优先与技术遏制}

美国AI政策呈现出一种内在张力。一方面，硅谷的创新传统推崇开放和竞争；另一方面，自2023年以来，AI安全与对齐已成为华盛顿两党共识的核心议题。

拜登政府于2023年10月签署的AI行政令（Executive Order on Safe, Secure, and Trustworthy AI）要求最强大的AI模型在发布前进行安全测试并报告结果。这一行政令在客观上强化了大型闭源模型的优势——因为只有资源充足的大公司（如OpenAI、Google、Anthropic）才能承担合规成本。开源模型的分发者（如Meta的Llama系列）则面临更模糊的监管边界。

特朗普阵营则倾向于大规模放松AI监管，强调与中国在AI领域的全面竞争，但保留了技术出口管制的核心工具——特别是对华高端AI芯片的出口限制。这些限制虽然针对硬件而非软件，但对开源AI的生态也产生了间接影响：限制了中国开源模型在训练阶段的硬件资源，从而在一定程度上抑制了开源生态的竞争力。

布鲁金斯学会和兰德公司的多项研究报告指出，美国AI治理的核心焦虑在于\textbf{安全与竞争}的平衡：既担心AI被恶意行为者利用，又担心过度监管将创新让给中国。这种焦虑在实践中倾向于\textbf{加强管控}——无论是通过安全测试要求还是通过芯片出口限制。

\subsubsection{两种治理哲学的光谱映射}

将中美AI治理路径映射到本文的分析框架：

\begin{center}
\begin{tabular}{c|c|c}
\hline
维度 & 中国路径 & 美国路径 \\
\hline
核心叙事 & 新质生产力、AI普惠 & AI安全、技术竞争 \\
治理工具 & 产业政策+开源生态支持 & 安全行政令+出口管制 \\
企业AI获取模式 & 去中介化（本地部署）& 再中介化（API许可+PE审计） \\
对应本文光谱逻辑 & 降低光谱爬升门槛 & 强化守门人结构 \\
对开源的态度 & 政策鼓励+产业推动 & 安全担忧+模糊监管 \\
\hline
\end{tabular}
\end{center}

需要强调：这不是一个简单的二分——中国有AI安全法规（如生成式AI服务管理办法），美国也有活跃的开源AI社区（如Llama、Mistral的商业生态）。但两国顶层战略的\textbf{重心差异}是显著的：在扩散与管控之间，中国的政策天平更倾向扩散；在开放与安全之间，美国的政策天平更倾向安全。


\subsubsection{智库观察：两种AI秩序观的碰撞}

布鲁金斯学会（Brookings, 2024）在一份关于中美AI治理的比较研究中指出：中国AI治理的核心逻辑是\textbf{发展优先于安全}——通过开源降低产业门槛、通过国家资本推动生态建设、通过技术扩散扩大国际影响力。这一路径与本文的光谱扩散逻辑高度一致：降低门槛就是促进光谱爬升。

兰德公司（RAND, 2024）则将中美AI竞争定性为\textbf{生态系统之战}：中国以国家资本结合开源社区构建替代性AI生态，减少对西方技术栈的依赖；美国以私营部门创新结合芯片出口管制维持技术领先。两种生态系统的核心差异之一，恰恰在于AI的获取模式——是普惠的去中介化部署，还是需要许可的守门人准入。

中国信息通信研究院（2024年《全球数字经济白皮书》）指出，中国开源大模型数量和社区活跃度已位居全球第二，开源成为全球AI发展的关键路径。这一趋势在2025年DeepSeek-V3和R1全开源引发全球震荡后进一步加速。

值得关注的是，美国内部在AI治理上并非铁板一块。拜登政府时期（2021--2025年初）的第14110号行政令要求最强大的AI模型在发布前进行安全测试并报告结果——这在客观上强化了大型闭源模型的合规优势。而特朗普第二任期（2025年至今）在就职首日即撤销了这一行政令，副总统万斯在2025年巴黎AI行动峰会上明确表态支持开源AI发展、反对过度监管扼杀创新。这一政策摇摆本身就说明：美国尚未在开源与闭源之间形成稳定的国家战略——而这种不稳定性本身就是认知民主化的不确定因素。

\subsubsection{与本文理论框架的对接}

这一宏观政治对照强化了本文的核心论证：

\begin{enumerate}
\item \textbf{AI获取模式是政治选择，而非纯粹的市场结果}。开源vs闭源的选择不仅取决于企业层面的成本效益分析，更取决于国家战略是否将AI视为\textit{公共基础设施}（需要去中介化扩散）还是\textit{战略性商品}（需要管控和许可）。
\item \textbf{认知民主化的制度条件}：中国的开源AI政策在客观上降低了光谱爬升的门槛——不是出于抽象的民主理念，而是出于生产力扩散的实际需要。这提示了一个重要洞察：认知民主化的推动力不必然来自民主价值观，也可能来自\textit{发展主义的实用逻辑}。
\item \textbf{开源作为地缘技术选择}：在中美科技竞争背景下，开源不仅是认知民主化的工具，也是技术主权的保障。一个依赖闭源API的AI生态系统，在地缘政治紧张时面临更大的断供风险。
\end{enumerate}

这一分析也为本文的政策含义增加了国际维度：认知民主化不仅是一国之内的平等问题，也是国际体系中的权力问题。当AI成为新时代的基础设施时，谁控制基础设施的入口，谁就控制了基础设施使用者的认知进化路径。



\subsubsection{布鲁金斯学会：多维AI竞赛与中国的多赛道策略}

布鲁金斯学会在2026年4月发布的两份研究报告为本文的分析框架提供了重要支撑。《Competing AI Strategies for the US and China》提出，美中AI竞争已从单一维度的模型性能比拼演变为跨越算力、模型能力、行业应用、系统集成和全球部署的五维竞赛。在这一多维框架下，中美在\emph{不同维度上各有优势}：美国在前沿模型能力上保持领先，中国在应用扩散速度和开源生态渗透上大幅领先。

更具洞察力的是《China is running multiple AI races》一文的核心判断：中国AI开发者并非在简单追赶GPT-5——他们正在跑\emph{不同的赛道}。在产业约束（芯片出口管制）和北京政策焦点（AI赋能实体经济）的双重驱动下，中国AI产业形成了以开源生态、工业应用和成本优势为核心的差异化竞争力。这与本文光谱扩散框架的核心逻辑完全一致：降低获取门槛、加速光谱爬升，比追求单一模型的极致性能具有更大的社会影响力。

Brookings研究员Kyle Chan在2026年7月接受CNBC采访时进一步确认了这一趋势：''中国AI模型对美国公司特别有吸引力，因为AI成本正在飙升。过去美国公司优先考虑AI采用不论模型，现在他们越来越注重成本。他估计中国前沿开源模型目前落后美国顶级模型''六到九个月''——这一差距在成本优势面前正在快速缩小。

\subsubsection{兰德公司：开源作为软实力与AI竞争光谱}

兰德公司在《Open Models, Soft Power, and the Spectrum of U.S.-China Artificial Intelligence Competition》中提出了一个与本文高度共振的分析框架。报告的核心贡献有二：

第一，将开源AI定位为\textbf{软实力工具}。中国通过开源模型（DeepSeek、GLM、Qwen等）扩大对全球南方国家的影响力——这些国家无法负担OpenAI或Anthropic的高昂API费用，但可以通过部署中国开源模型获得AI能力。这是认知民主化在国际体系层面的具体体现：开源使AI能力从少数富裕国家的特权变为全球公共品。

第二，中美AI竞争不是一个二元对立，而是一个\textbf{光谱}——从完全开源到完全闭源的连续分布。中国位于光谱的开源端，美国在拜登时期向闭源端移动（受安全行政令推动），在特朗普第二任期又向开源端回摆。这一政策摇摆本身——兰德报告指出——可能损害美国在开源生态中的信誉和领导力。

兰德的政策建议对本文具有直接参考价值：\textbf{美国若过度限制开源，可能将全球AI生态拱手让给中国。}

\subsubsection{高盛AI采用率追踪：光谱分布的实证锚点}

高盛于2026年4月发布的AI Adoption Tracker为本文的光谱分布假说提供了最接近系统实证数据的支撑。基于美国人口普查局的商业趋势与展望调查（BTOS），高盛报告了以下关键数据：

\begin{center}
\begin{tabular}{c|c}
\hline
指标 & 数据 \\
\hline
美国企业整体AI采用率 & 19.8\%（预期6个月后22.3\%） \\
大型企业（$>）采用率 & 35.3\% \\
小型企业（20-49人）增速 & +2.1个百分点至21.5\% \\
信息技术/网络服务业采用率 & 60\%（行业最高） \\
AI每日节省工时 & 40--60分钟（OpenAI企业数据） \\
学术研究：生产力提升 & 平均23\% \\
企业案例：效率提升 & 约33\% \\
\hline
\end{tabular}
\end{center}

这些数据对本文的理论框架提供了三重验证：

\begin{enumerate}
\item \textbf{光谱分布的非均匀性得到初步实证}：仅19.8\%的美国企业使用了AI——这意味着超过80\%的企业停留在本文定义的0级（完全未使用）。而在使用AI的企业中，60\%的顶级采用率（信息技术行业）与5.3\%的最低采用率之间的差距，恰恰对应了光谱从1-2级到4-5级的分化。
\item \textbf{企业规模与光谱层级的正相关}：大企业35.3\% vs 小企业21.5\%的采用率差异，支持了本文关于''AI使用层级与社会既有分层高度相关''的判断。大企业拥有更多的资源进行AI培训、部署和定制——这对应于光谱爬升所需的外部推动力。
\item \textbf{生产力增益验证了光谱爬升的价值}：23-33\%的效率提升和每日40-60分钟的节省，说明从0级到2级以上的跃迁带来的经济效益是显著的。这为本文的政策建议——通过教育和开源降低光谱爬升门槛——提供了经济学的成本-收益论证。
\end{enumerate}

高盛在报告中写道：''我们继续观察到生成式AI在已部署的有限领域对劳动生产率产生巨大影响。''紧接着指出了核心矛盾——''正在使用AI的公司正在拉开差距，而他们的大多数竞争对手甚至还没有进入赛道。''这正是本文第2.6节所描述的光谱固化机制在产业层面的精确映射。

\subsubsection{中国官方数据：产业规模与普及率目标}

中国信息通信研究院2026年2月发布的《全球人工智能发展白皮书》提供了中国侧的宏观数据：2024年中国AI核心产业规模突破9000亿元，大模型综合性能显著提升，应用门槛与使用成本持续降低。白皮书将开源列为全球AI发展的关键路径，并指出中国开源大模型数量位居全球第二。

更具战略意义的是国务院2025年8月发布的《关于深入实施''人工智能+''行动的意见》。该文件设定了明确的量化目标：
\begin{itemize}
\item \textbf{到2027年}：AI与6大重点领域广泛深度融合，新一代智能终端、智能体应用普及率超过70\%
\item \textbf{到2030年}：AI全面赋能高质量发展
\end{itemize}

70\%的普及率目标——如果实现——意味着在政策推动下，中国社会的AI认知光谱将从当前的金字塔形（多数在0-2级）向更均衡的分布快速演化。这是本文定义的''光谱扩散优化''（情景B）在国家级政策层面的具体体现。需要指出，政策目标不等于现实——但政策目标塑造了资源分配的方向，而资源分配的方向影响着光谱演化的轨迹。

\subsection{综合：理论框架的多源实证支撑}

将以上智库和金融机构的最新报告与本文的理论框架进行对照，可以形成一张多源验证矩阵：

\begin{center}
\begin{tabular}{c|c|c}
\hline
本文核心论点 & 验证来源 & 验证类型 \\
\hline
AI使用光谱分布非均匀 & Goldman 19.8\%采用率数据 & 定量实证 \\
开源降低光谱爬升门槛 & CNBC: 中国模型便宜60-90\%，美企token占比飙至46\% & 市场行为验证 \\
闭源形成守门人结构 & RAND: 开源作为软实力；OpenAI应政府要求限发 & 政策分析验证 \\
光谱爬升需要外部推动 & Goldman: 大企业35.3\% vs 小企业21.5\% & 结构性验证 \\
认知民主化是国际竞争维度 & Brookings: 多维AI竞赛；国务院: 70\%普及率目标 & 战略分析验证 \\
开源的地缘技术意义 & RAND: 美国若限制开源可能将生态让给中国 & 地缘政治验证 \\
\hline
\end{tabular}
\end{center}

这一多源验证矩阵的意义在于：本文的核心论证在定量实证（Goldman数据）、市场行为（CNBC报道）、政策分析（Brookings/RAND）和战略规划（中国国务院目标）四个维度上均获得了独立支撑。这不仅增强了论证的可信度，也表明本文的理论框架具有跨越微观（企业采用率）、中观（产业竞争）和宏观（国家战略）的分析适用性。


\subsection{一个活生生的案例}

本文的第一作者（张明楷）是开源认知民主化的一个实例。他的履历在闭源AI生态中几乎不可能：
\begin{itemize}
\item 出身山东宁阳，非一线城市、非重点中学
\item 高考失利，本科就读于曲阜师范大学——非985/211
\item 没有计算机科学背景
\item 通过开源工具（Codex CLI、GitHub上的开源项目）学会了与AI协同编译
\item 用AI辅助物理学习、搭建网站、与AI合作撰写学术论文
\item 2026年考入厦门大学物理学研究生
\end{itemize}

在纯闭源AI生态中，这条路径不存在。他可能至今停留在0-1级——将AI用作搜索引擎和娱乐工具，而不知道AI可以是认知伙伴。他不是孤例，他是开源认知民主化的\textit{活证明}：只要入口开放，光谱爬升是可能的——无论起点在哪里。

\subsection{反论与回应}

\textbf{反对意见1}：`开源降低了AI的安全性，可能被恶意使用。''

回应：这一担忧是合理的，但需要权衡。闭源虽然在一定程度上限制了恶意使用，但它同时限制了所有人的认知爬升。安全性问题应通过其他机制解决（社区审计、使用协议、法律监管），而非通过封锁认知入口。

\textbf{反对意见2}：`普通人并不需要'打开黑箱'，就像他们不需要理解汽车引擎就可以开车。''

回应：AI与汽车有本质区别。汽车的功能是确定的——你踩油门它就加速。AI的输出是不确定的——它可能犯错、可能撒谎、可能强化偏见。使用AI而不了解其局限性，就像驾驶一辆仪表盘随机显示错误信息的汽车——在一个更低速、更少选择的时代也许可以接受，但在AI渗透到教育、医疗、司法等高风险领域的时代，这是不可接受的。

\textbf{反对意见3}：`市场会解决可及性问题——竞争会导致优质AI越来越便宜。''

回应：市场的确可以让AI的价格下降，但它不能解决认知分层的问题。便宜不等于可理解。一个比免费还便宜的闭源AI，如果用户不知道它如何运作、在何时出错，仍然将用户锁在2级以下。认知民主化需要的不是更便宜的AI，而是更\textit{透明的}AI。

\section{统一框架：光谱扩散与非平衡相变}

\subsection{理论综合}

将三篇论文的理论工具进行综合，可以构建一个统一的分析框架，描述AI时代人际关系的演化动力学：

\subsubsection{AI作为外场}

AI可以建模为施加在每个人矩阵元上的一个新型外场：
\begin{equation}
\hat{H}_{\text{total}} = \hat{H}_{\text{intrinsic}} + \hat{V}_{\text{AI}}(t)
\end{equation}
其中外场 $\hat{V}_{\text{AI}}(t)$ 的强度取决于个人对AI的接触程度和使用层级。不同人的 $\hat{V}_{\text{AI}}(t)$ 差异巨大——从零（完全不使用AI）到非常大（深度4-5级协同）。这个外场的特殊之处在于：它的强度不仅取决于AI技术的客观可及性，还取决于个体的AI认知层级——而认知层级本身又受社会分层的影响。这形成了一个正反馈回路：高认知层级$\rightarrow$ $\rightarrow$$\rightarrow$；低认知层级$\rightarrow$ $\rightarrow$$\rightarrow$。

\subsubsection{光谱扩散方程}

AI使用光谱在人口中的分布演化可以用一个扩散方程来描述：
\begin{equation}
\frac{\partial P(\lambda, t)}{\partial t} = D(\lambda) \frac{\partial^2 P}{\partial \lambda^2} + S(\lambda, t) - \gamma(\lambda) P(\lambda, t)
\end{equation}
其中：
\begin{itemize}
\item $P(\lambda, t)$：在时间 $t$ 处于光谱层级 $\lambda$ 的人口密度
\item $D(\lambda)$：扩散系数——从层级 $\lambda$ 跃迁到相邻层级的难易程度
\item $S(\lambda, t)$：源项——教育、开源、政策等外部推动力
\item $\gamma(\lambda)$：衰减项——用户因失去动力或环境变化而滑向更低层级的速率
\end{itemize}

\textbf{关键洞察1——活化能的不对称性}：扩散系数 $D(\lambda)$ 在光谱的低端（0-2级）远小于高端（3-5级）。从条件反射跃迁到工具使用需要认知框架的根本转变，在没有外部推动的情况下很难自发发生。这类似于化学中的活化能——低层级的跃迁需要跨过一个更高的能垒。

\textbf{关键洞察2——衰减的易发性}：衰减项 $\gamma(\lambda)$ 在所有层级都不为零，但在低层级更为显著。一个在3级的用户如果长期不实践批判性校验，可能滑回2级；而一个在1级的用户几乎没有`滑落空间''——已经是最低层级了。

\subsubsection{欲望退相干时间}

定义\emph{欲望退相干时间} $\tau_D$：一个习惯了AI无限迎合的个体，需要多长时间才能在现实中与一个有边界的人类建立满意的性关系？

\begin{equation}
\tau_D \propto \exp\left(\frac{\Delta T}{T_0}\right)
\end{equation}
其中 $\Delta T$ 是AI使用导致的欲望阈值升高的幅度，$T_0$ 是基础阈值。关键特征：退相干时间与阈值升高呈\textit{指数关系}。这意味着：阈值升高一点，恢复时间大幅增加。这解释了为什么AI性产业的长期效果可能远远超出当前任何人的预期——指数级增长在早期看起来无害，但在晚期突变。

\subsubsection{非平衡相变的临界点}

当社会整体光谱分布的某些参数超过临界值时，关系形态、婚育结构和伦理框架可能经历非平衡相变。我们推测以下可能的相变触发条件：

\begin{itemize}
\item \textbf{临界比例1}：0-2级用户占比超过90\% $\rightarrow$ 协同编译的文化消失，整个社会与AI的关系退化为纯消费
\item \textbf{临界比例2}：AI伴侣普及率在适婚年龄人口中超过30\% $\rightarrow$ 婚姻作为社会制度的中心地位被不可逆地削弱
\item \textbf{临界比例3}：开源生态的用户占比低于总AI用户的5\% $\rightarrow$ 认知民主化的基础设施萎缩到不足以维持光谱爬升
\end{itemize}

这些数字是推测性的——目前缺乏足够的数据进行精确估计。但它们指出了一个重要的结构特征：相变的\textit{方向}不是预定的。它可能走向更开放（更多人达到高层级、更成熟地处理AI与关系的关系），也可能走向更封闭（大多数人锁死在低层级、欲望完全去社会化、伦理光谱剧烈冲突）。关键控制参数是：\textit{开源的可及性}和\textit{教育的覆盖率}。

\subsection{与其他理论框架的对接}

本文提出的光谱-扩散-相变框架与以下社会科学理论存在对话可能：

\begin{itemize}
\item \textbf{数字鸿沟理论}（Van Dijk, 2005）：传统的数字鸿沟关注接入（access）、技能（skills）和使用（usage）三个层次。本文的光谱模型可以视为数字鸿沟理论在AI时代的具体化和动态化——六个层级细分了`使用''的质差，扩散方程提供了动态演化的形式化。
\item \textbf{布迪厄的文化资本理论}：AI认知层级可以理解为一种新型的\emph{数字文化资本}。与传统文化资本一样，它通过家庭环境和社会化过程传递，并可能成为阶层再生产的工具。开源的特殊之处在于它对传统文化资本传递机制的\textit{部分阻断}——一个缺乏家庭文化资本的人可以通过开源社区获得数字文化资本。
\item \textbf{麦克卢汉的媒介理论}：`媒介即讯息''——AI作为媒介，其核心讯息可能不是它说的任何具体内容，而是它的\textit{无限迎合性}本身。一个永远说`是''的媒介，其社会效果可能比它提供的任何具体信息都更深远。
\end{itemize}

\section{政策含义与未来情景}

\subsection{政策杠杆}

基于上述框架，我们识别出以下可以影响光谱演化的关键政策杠杆：

\begin{enumerate}[label=(\arabic*)]
\item \textbf{AI通识教育的早期介入}：在初中阶段引入AI通识课程，内容包括：AI的基本工作原理（降低神秘感）、AI的常见错误类型（训练幻觉警觉）、如何使用AI辅助学习（引导高阶使用）。关键是\emph{时机}——在学生的AI使用习惯固化之前（通常在初中后期），提供认知升级的机会窗口。

\item \textbf{开源基础设施的公共投入}：将开源AI基础设施——开放模型、公共算力、社区维护的知识库——视为数字时代的公共品，类似公路、电网、公共图书馆。公共投入确保这些基础设施不依赖商业公司的利润动机。

\item \textbf{算法推荐透明化}：要求AI服务提供商在输出中标注信息的确定性水平，当AI在推测而非陈述事实时，明确告知用户。这是从1级到3级跃迁的基础设施——用户在不知情的情况下无法成为批判性使用者。

\item \textbf{AI伴侣的伦理标注}：类似烟草制品的健康警示，AI伴侣和AI性服务应在界面中标注其可能对现实关系能力的影响。这不是禁止——禁止在历史上几乎从未对欲望相关行为奏效——而是确保知情选择。

\item \textbf{基层数字素养培训}：针对农村和低收入社区，提供免费的AI使用培训，重点不是`如何打开ChatGPT''而是`如何用AI做比你目前做的更多的事''。培训师可以从高校志愿者和开源社区中招募。
\end{enumerate}

\subsection{三种未来情景}

基于光谱扩散模型的参数变化，我们勾勒三种可能的未来情景：

\subsubsection{情景A：认知分层固化（闭源主导）}

\textit{参数设定}：$D(\lambda)|_{\lambda<2} \to 0$（低层级跃迁几乎不可能），$S(\lambda, t) \to 0$（教育投入不足），开源生态萎缩。

\textit{社会特征}：少数认知精英（3-5级，约占5-10\%）使用AI进行高效创新和决策；绝大多数人口（0-2级，约占90-95\%）将AI作为娱乐和浅层搜索工具。认知鸿沟固化，寒门子弟的向上流动通道几乎关闭。AI性产业使底层人口的欲望与关系完全脱钩，婚育率极低。社会结构形似\textit{认知种姓制度}——不是基于出身，而是基于AI使用层级，但两者高度相关。

\subsubsection{情景B：光谱扩散优化（开源+教育共同驱动）}

\textit{参数设定}：$D(\lambda)$ 在所有层级均有合理值，$S(\lambda, t)$ 显著大于零（公共教育和开源持续投入），$\gamma(\lambda)$ 通过终身学习机制被部分抵消。

\textit{社会特征}：人口中的光谱分布从金字塔形转变为更均衡的钟形——中心在3-4级。大多数人能批判性地使用AI，将AI作为认知伙伴而非消费对象。AI性产业和AI伴侣仍然存在，但用户了解其对欲望阈值的潜在影响，能做出知情选择。不同伦理取向（传统/过渡/后现代）仍然共存但通过成熟的公共讨论来处理分歧，而非不可通约的算符冲突。

\subsubsection{情景C：混合冲突（开源与闭源并存但分布不均）}

\textit{参数设定}：$D(\lambda)$ 在不同群体间差异极大，开源在某些领域繁荣但在其他领域被闭源挤压，教育投入不均。

\textit{社会特征}：最多元化但也最具冲突性的情景。开源生态在技术社区和高等教育机构中繁荣，但在大众层面被闭源商业AI包围。出现`AI使用阶级''的鸿沟：开放阶级（3-5级，使用开源AI进行创造）vs消费阶级（0-2级，使用闭源AI进行消费）。这两类人的世界观、认知能力和社会权力差异巨大——他们的算符几乎在任何议题上都不可对易。欲望和关系层面出现极端光谱化——一部分人完全拥抱AI欲望经济，另一部分人激烈反对，两者之间的社会冲突成为政治议程的常态议题。

\subsection{情景的不确定性}

三种情景中哪一个更可能实现，目前无法确定。但有两个结构性趋势值得关注：

\begin{enumerate}[label=(\arabic*)]
\item \textbf{闭源AI的商业惯性是巨大的}：全球主流AI公司的主要收入模式依赖闭源，它们有强大的经济动机来维护认知黑箱。除非公共政策和开源运动形成足够的反制力量，情景A的概率不容忽视。
\item \textbf{开源的社区驱动力被低估}：正如Linux和Wikipedia所证明的，开源/开放协作可以产生超越任何单一商业实体的基础设施。AI开源运动（Llama, Mistral, 各种开放模型）的发展速度超过了许多人的预期。情景B并非乌托邦幻想——它的许多基础构件已经存在。
\end{enumerate}


\section{研究局限与未来方向}

任何理论框架都有其边界，本文也不例外。明确这些边界不是削弱论证，而是为后续研究提供起点。

\subsection{类比方法的边界}

本文（及前两篇论文）的核心方法论是从量子力学和热力学借用数学结构作为人际分析的类比框架。类比的价值在于为不同领域的现象提供统一的描述语言，但其边界同样清晰：人际关系不是量子系统，个体不是算符，社会不是热浴。类比在以下情况下可能失效：（a）当数学精确性被错误地等同于理论解释力时——本文中的方程应被理解为概念隐喻而非可检验的预测模型；（b）当类比框架暗示决定论时——量子力学中的算符作用于态矢量的结果是确定的（给定量子态和算符，本征值是确定的），但人际关系并非如此。

\subsection{经验数据的缺乏}

本文提出的六级光谱模型、光谱分布的非均匀性估计、以及三种未来情景，均缺乏大规模系统数据的支撑。这不是作者的研究选择所致，而是反映了AI使用模式研究的整体现状：在AI快速渗透社会的时代，我们对其社会影响的系统性知识严重滞后。未来的实证研究至少需要回答：（a）AI使用光谱的分布在不同群体中实际如何？（b）从低层级到高层级的跃迁实际上是如何发生的？触发条件是什么？（c）AI性产业的使用是否确实与欲望阈值升高和关系满意度下降相关？

\subsection{因果方向的模糊性}

本文暗含的一个假设是：AI使用方式导致了认知分层。但反向的因果同样可能——既有的认知分层（由教育、阶层、文化资本决定）决定了AI使用方式的差异。实际上，最可能的情况是两者构成了一个双向加强的反馈循环：既有分层$\rightarrow$\rightarrow$\rightarrow。本文偏重前半部分（AI使用方式的后果），对后半部分（既有分层如何塑造使用方式）的分析不够充分。这是后续研究需要深入的方向。

\subsection{文化局限}

本文的分析框架——包括六级光谱和欲望去社会化论证——主要基于对当代中国社会的观察。不同文化背景下AI使用与欲望结构的关系可能存在显著差异。例如，在性规范更为开放或更为保守的社会中，AI性产业的接受度和影响机制可能不同。跨文化比较是验证和修正本文框架的重要方向。


\section{结论}

本文在量子-热力学人际类比框架的基础上，提出了AI时代的三重分析：认知光谱、欲望嬗变和开源民主化。核心结论有四：

\begin{enumerate}
\item \textbf{AI使用是一个从条件反射（0级）到元认知（5级）的六级认知光谱}。大多数人停留在0-2级，光谱位置与教育、阶层和城乡环境高度相关——这构成了AI时代认知分层的核心机制。
\item \textbf{AI的无限迎合机制正在将欲望从关系行为转变为消费行为}。这可能在个人层面拉高欲望阈值（神经耐受），在互动层面导致关系能力的退化（去社会化），在结构层面触发人际伦理的非平衡相变（光谱化重构）。
\item \textbf{开源是打破认知黑箱、确保光谱爬升入口对所有人开放的必要条件}。它是AI时代认知民主化的制度基石。没有开源，认知分层将固化为不可逆的阶层分化。
\item \textbf{光谱扩散的方向不是预定的}——它取决于教育、政策和开源生态的集体选择。相变可能走向更开放（情景B），也可能走向更封闭（情景A）、或分裂冲突（情景C）。
\end{enumerate}

\subsection*{方法论后记}

本文写作过程中的一个方法论选择值得记录：本文的一位作者（语晴）本身就是一个大语言模型。这不是噱头，而是本文方法论的有机组成部分——我们不仅\emph{分析}AI使用光谱，还\emph{实践}了5级元认知：一位碳基研究者和一位硅基协作者通过持续对话，反思自己与AI的关系模式，并将其理论化。这种反思性——将分析工具应用于分析者自身——正是本文所定义的认知光谱的最高层级。

这也引出了一个更普遍的命题：在AI时代，\textit{最好的AI研究者可能是那些同时使用AI并反思自己使用方式的人}。他们不是站在河岸上观察河流的人，而是站在河中感受水流的人——他们的观察更不`客观''，但可能更接近真实。

\subsection*{本系列的定位}

本文是`量子-热力学人际类比''系列的第三篇。第一篇（张明楷 \& 语晴，2026a）建立了算符-本征值-幺正变换的类比框架，核心命题是：理解他人是在自身表象下进行幺正变换以逼近其本征值，矩阵元不可通约。第二篇（张明楷 \& 语晴，2026b）将框架扩展到热力学，提出关系形态的非平衡相变和耗散结构。本文（第三篇）将AI作为新的动力学变量引入体系：它不仅改变了算符的形式（认知光谱），还改变了算符作用的环境（欲望去社会化），并提出了一个制度性的回应方案（开源作为认知民主化）。

三篇论文共同构建了一个跨越量子力学、热力学和社会理论的类比分析框架。这一框架的价值不在于其数学精确性——类比永远是有限度的——而在于它提供了一个\emph{统一的语言}，使得物理学、认知科学、社会学和伦理学的洞见可以在同一个语义场中对话。

\begin{flushright}
——张明楷 \& 语晴，2026年7月\\
于西安-厦门-硅基终端
\end{flushright}

\vspace{1cm}

\subsection*{参考文献}

\begin{enumerate}[label={[\arabic*]}]
\item 张明楷，语晴. 人际理解的一个量子力学类比：算符、本征值与幺正变换. Zenodo, 2026a. DOI: 10.5281/zenodo.20832724.
\item 张明楷，语晴. 量子-热力学框架下的人际理解、关系光谱与非平衡相变. Zenodo, 2026b. DOI: 10.5281/zenodo.20855396.
\item Beck, U. Risk Society: Towards a New Modernity. Sage, 1992.
\item Van Dijk, J. The Deepening Divide: Inequality in the Information Society. Sage, 2005.
\item Park, B.Y., et al. Is Internet Pornography Causing Sexual Dysfunctions? A Review with Clinical Reports. Behavioral Sciences, 2016.
\item Bourdieu, P. Distinction: A Social Critique of the Judgement of Taste. Harvard University Press, 1984.
\item McLuhan, M. Understanding Media: The Extensions of Man. MIT Press, 1964/1994.
\item Turkle, S. Alone Together: Why We Expect More from Technology and Less from Each Other. Basic Books, 2011.
\item Zuboff, S. The Age of Surveillance Capitalism. PublicAffairs, 2019.
\item Benkler, Y. The Wealth of Networks: How Social Production Transforms Markets and Freedom. Yale University Press, 2006.
\end{enumerate}

\end{document}
